%!TEX root = ../template.tex
%%%%%%%%%%%%%%%%%%%%%%%%%%%%%%%%%%%%%%%%%%%%%%%%%%%%%%%%%%%%%%%%%%%%
%% abstract-pt.tex
%% NOVA thesis document file
%%
%% Abstract in Portuguese
%%%%%%%%%%%%%%%%%%%%%%%%%%%%%%%%%%%%%%%%%%%%%%%%%%%%%%%%%%%%%%%%%%%%

\ExplSyntaxOn
\cs_if_exist:NT \TitleOnHeader
  { \TitleOnHeader }
\ExplSyntaxOff

\typeout{NT FILE abstract-pt.tex}%

Uma das principais ameaças a sistemas de anonimato como o Tor é o denominado Sybil attack, em que o atacante personifica várias identidades para uma só entidade. Para resolver esta ameaça, eu proponho o Dittor, um esquema de credenciais anónimo e resistente a Sybil attacks. 

O Dittor faz recurso a credenciais aprovadas por entidades externas para validação da identidade única de um operador de nó de rede do Tor, de forma anónima. Este sistema de credenciais não permite que haja mais do que uma identidade para um só operador de nó, mitigando assim possíveis Sybil attacks. Enquanto tentativas anteriores procuravam a identificação de nós sybil na rede Tor para mitigação, o Dittor resolve o problema de forma fundamental, associando credenciais anónimas, que validam a identidade do sujeito, a um pseudónimo, que é então usado como identificador na rede do Tor para agrupar os nós de uma entidade numa só ligação, sem revelar a identidade do mesmo. Para atingir este fim, mecanismos como Sybil-resilient Anonymous Signatures e Zero-Knowledge Proofs são usados para a construção de um sistema de Self-Sovereign Identity, culminando num sistema de Sybil-resilient Anonymous Signatures with Stateless Issuing.

Eu proponho um método de trabalho de nove meses, envolvendo fases de pesquisa e desenvolvimento, com um gráfico temporal para simplificação dos objetivos.


% E agora vamos fazer um teste com uma quebra de linha no hífen a ver se a \LaTeX\ duplica o hífen na linha seguinte se usarmos \verb+"-+… em vez de \verb+-+.
%
% zzzz zzz zzzz zzz zzzz zzz zzzz zzz zzzz zzz zzzz zzz zzzz zzz zzzz zzz zzzz comentar"-lhe zzz zzzz zzz zzzz
%
% Sim!  Funciona! :)

% Palavras-chave do resumo em Português
% \begin{keywords}
% Palavra-chave 1, Palavra-chave 2, Palavra-chave 3, Palavra-chave 4
% \end{keywords}
\keywords{
  Tor \and
  Sybil \and
  Identificadores Descentralizados \and
  Autonomia de Identidade
}
% to add an extra black line
